\begin{lemma}
\label{lemma-tensor-product}
Let $M, N$ be $R$-modules. Then there exists a pair $(T, g)$
where $T$ is an $R$-module, and
$g : M \times N \to T$ an $R$-bilinear
mapping, with the following universal property:
For any $R$-module $P$ and any $R$-bilinear mapping
$f : M \times N \to P$, there
exists a unique $R$-linear
mapping $\tilde{f} : T \to P$ such that $f = \tilde{f} \circ g$.
In other words, the following diagram commutes:
$$
\xymatrix{
M \times N \ar[rr]^f \ar[dr]_g & & P\\
& T \ar[ur]_{\tilde f}
}
$$
Moreover, if $(T, g)$ and $(T', g')$
are two pairs with this property, then there
exists a unique isomorphism
$j : T \to T'$ such that $j\circ g = g'$.
\end{lemma}

\noindent
The $R$-module $T$ which satisfies the above universal property is called
the  {\it tensor product} of $R$-modules $M$ and $N$, denoted as
$M \otimes_R N$.

\begin{proof}
We first prove the existence of such $R$-module $T$.
Let $M, N$ be $R$-modules.
Let $T$ be the quotient module
$P/Q$, where $P$ is the free $R$-module $R^{(M \times N)}$ and $Q$ is the
$R$-module generated by all elements of
the following types: ($x\in M, y\in N$)
\begin{align*}
(x + x', y) - (x, y) - (x', y), \\
(x, y + y') - (x, y) - (x, y'), \\
(ax, y) - a(x, y), \\
(x, ay) - a(x, y)
\end{align*}
Let $\pi : M \times N \to T$ denote the natural map.
This map is $R$-bilinear, as
implied by the above relations
when we check the bilinearity conditions. Denote the image
$\pi(x, y) = x \otimes
y$, then these elements generate
$T$. Now let $f : M \times N \to P$ be an $R$-bilinear map,
then we can define
$f' : T \to P$ by extending the mapping
$f'(x \otimes y) = f(x, y)$. Clearly $f = f'\circ \pi$. Moreover, $f'$ is
uniquely determined by the value on the
generating sets $\{x \otimes y : x\in M, y\in N\}$.
Suppose there is another pair $(T', g')$ satisfying the same properties.
Then there is a unique $j : T \to T'$ and
also $j' : T' \to T$ such that $g' = j\circ g$, $g = j'\circ g'$.
But then both the maps $(j\circ j') \circ g$ and $g$
satisfies the universal properties, so by uniqueness they are equal,
and hence $j'\circ j$ is identity on $T$.
Similarly $(j'\circ j) \circ g' = g'$ and $j\circ j'$ is identity on $T'$.
So $j$ is an isomorphism.
\end{proof}